Tool-augmented data agents can execute a tool successfully yet reach a wrong conclusion from a corrupted response. We introduce ToxicBench, a diagnostic benchmark that pairs clean and poisoned tool observations while keeping source data fixed. Its tasks cover numerical, label, schema, and retrieval errors, with trajectory metrics separating checking from final-answer adoption. Across three adapters in the 118-task GPT evaluation, poisoned task success falls by 26--39 percentage points. Two-route comparisons show the benefit of an ordinary retry under one-shot poisoning, while repeated poisoning reveals continued adoption of wrong answers after checking. Controlled experiments on three public tables further isolate the effect of supplied evidence on recovery. A post-freeze human evaluation of 200 trajectories yields 96\% task-success scoring agreement and corroborates both the retry benefit over Base on the audited tasks and adoption after checking. We release trajectories, versioned scoring, and reference and delivery audits. These results establish evidence availability and final-answer selection as complementary dimensions for evaluating tool-augmented data agents.
